\mtexe{3.24}
\begin{proof}
	Equivalently, $P$ is totally ramified in $L$ if there is a unique prime $Q$ of $S$ lying over $P$ with inertial degree 1. Now, $Q \cap M$ is the unique prime of $M$ lying over $P$ since any prime of $M$ lying over $P$ must be contained in a prime of $S$ lying over $P$. We then also have:
	\[ 1 = f(Q \mid P) = f(Q \mid Q \cap M)f(Q \cap M \mid P) \]
	and so $f(Q \cap M \mid P) = 1$. I.e. $P$ is totally ramified in $M$. \\
	
	Notice that $K \subseteq L \cap L' \subseteq L$, so by the previous part, $P$ is totally ramified in $L \cap L'$. But since $K \subseteq L \cap L' \subseteq L'$, it is also unramified in $L \cap L'$. That is, if $U$ is a prime of $L'$ lying over $Q \cap L'$, then
	\[ 1 = e(U \mid P) = e(U \mid Q \cap L')e(Q \cap L' \mid P) \]
	and so
	\[ 1 = e(Q \cap L' \mid P) = [L \cap L' : K] \]
	so that $L \cap L' = K$ as claimed. \\
	
	As suggested, if we first take $m = p^r$, then we've shown in 2.34 that $p = u(1-\omega)^{\varphi(m)}$ for some unit $u$.
\end{proof}
